\documentclass[handout]{beamer}
\mode<presentation>
\usetheme[secheader]{Boadilla}
\usecolortheme{whale}

\usepackage[utf8]{inputenc}
\usepackage[T1]{fontenc}
\usepackage[indonesian]{babel}
\usepackage{lmodern}
\usepackage{graphicx}
\usepackage{bookmark}

\setbeamertemplate{footline}
{
	\leavevmode%
	\hbox{%
		\begin{beamercolorbox}[wd=.333333\paperwidth,ht=2.25ex,dp=1ex,center]{author in head/foot}%
			\usebeamerfont{author in head/foot}\insertshortauthor%
		\end{beamercolorbox}%
		\begin{beamercolorbox}[wd=.433333\paperwidth,ht=2.25ex,dp=1ex,center]{title in head/foot}%
			\usebeamerfont{title in head/foot}\insertshorttitle
		\end{beamercolorbox}%
		\begin{beamercolorbox}[wd=.233333\paperwidth,ht=2.25ex,dp=1ex,right]{date in head/foot}%
			\usebeamerfont{date in head/foot}\insertshortdate{}\hspace*{2em} \insertframenumber{} / \inserttotalframenumber\hspace*{2ex}
	\end{beamercolorbox}}%
	\vskip0pt%
}

\title[Bab 2 -- Pengantar Web]{Pemrograman Web}
\subtitle{Bab 2: Pengantar Pemrograman Web, Arsitektur Client-Server, Ruang Lingkup Kebutuhan Pengguna}
\author{Cecep Suwanda, S.Si., M.Kom.}
\institute{Universitas Bale Bandung}
\date{\today}

\begin{document}

% Slide Judul
\begin{frame}
	\titlepage
\end{frame}

% Outline dan CPMK
\begin{frame}{Outline Bab 2}
	\begin{block}{Sub-CPMK 1.1}
		Mengidentifikasi kebutuhan fungsional dan non-fungsional pengguna untuk solusi web.
	\end{block}
	\begin{itemize}
		\item Arsitektur Web: Client-Server dan Komunikasi HTTP
		\item Protokol HTTP: Metode, Status Code, dan Header
		\item Ruang Lingkup Kebutuhan Pengguna
		\item Teknologi Front-End: HTML, CSS, dan JavaScript
		\item Teknologi Back-End: PHP, MySQL, API
		\item Siklus Pengembangan Web (SDLC)
		\item Alat Pengembangan Web
		\item Standar Web dan Organisasi Terkait
		\item Peta Konsep: Hubungan Antar Teknologi Web
		\item Keamanan Web: Konsep Dasar
		\item Aksesibilitas dan Performa
	\end{itemize}
\end{frame}

% ========== SECTION 2-1: Arsitektur Web ==========
\begin{frame}{Arsitektur Web: Client-Server}
	\begin{itemize}
		\item \textbf{Client} (browser) mengirim permintaan; \textbf{server} menyimpan dan memproses data.
		\item Server mengembalikan respons: HTML, JSON, gambar, dll.
		\item \textbf{Analogi}: Client = pelanggan di restoran, Server = dapur, HTTP = pelayan.
		\item Komunikasi melalui \textbf{protokol HTTP}; transport \textbf{TCP}.
		\item \textbf{TCP three-way handshake}: SYN $\to$ SYN-ACK $\to$ ACK.
		\item Port \textbf{80} (HTTP), \textbf{443} (HTTPS).
	\end{itemize}
\end{frame}

\begin{frame}{DNS dan URL}
	\begin{itemize}
		\item \textbf{DNS}: menerjemahkan nama domain (\texttt{www.example.com}) ke alamat IP.
		\item \textbf{URL} (Uniform Resource Locator): alamat sumber daya di web.
	\end{itemize}
	\begin{block}{Komponen URL}
		\begin{itemize}
			\item \textbf{Skema}: \texttt{https} (protokol)
			\item \textbf{Host}: \texttt{www.example.com}
			\item \textbf{Port}: \texttt{:443} (opsional)
			\item \textbf{Path}: \texttt{/halaman/artikel}
			\item \textbf{Query}: \texttt{?id=42\&lang=id}
			\item \textbf{Fragment}: \texttt{\#bagian} (diproses di client)
		\end{itemize}
	\end{block}
\end{frame}

% ========== SECTION 2-2: Protokol HTTP ==========
\begin{frame}{Protokol HTTP: Metode dan Status Code}
	\begin{itemize}
		\item HTTP \textbf{stateless}; cookie/session untuk state.
		\item \textbf{Metode}: GET (ambil), POST (kirim), PUT (update penuh), PATCH (update sebagian), DELETE (hapus).
		\item \textbf{Idempotent}: GET, PUT, DELETE (hasil sama jika diulang).
	\end{itemize}
	\begin{block}{Status Code}
		\textbf{2xx} Sukses (200 OK, 201 Created) \\
		\textbf{3xx} Redirect (301, 302) \\
		\textbf{4xx} Error client (400, 401, 403, 404) \\
		\textbf{5xx} Error server (500, 503)
	\end{block}
\end{frame}

\begin{frame}{Header HTTP dan HTTPS}
	\begin{itemize}
		\item \textbf{Request header}: Host, User-Agent, Accept, Content-Type.
		\item \textbf{Response header}: Content-Type, Content-Length, Set-Cookie, Cache-Control.
		\item \textbf{HTTPS}: HTTP + enkripsi TLS/SSL; melindungi privasi dan integritas data.
		\item Membaca header: DevTools $\to$ tab \textbf{Network} $\to$ pilih request $\to$ Headers.
	\end{itemize}
\end{frame}

% ========== SECTION 2-1b: Kebutuhan Pengguna ==========
\begin{frame}{Ruang Lingkup Kebutuhan Pengguna}
	\begin{block}{Kebutuhan Fungsional}
		\textit{Apa} yang harus dilakukan sistem. Contoh: login, lihat produk, keranjang, checkout. \textbf{Wajib}.
	\end{block}
	\begin{block}{Kebutuhan Non-Fungsional}
		\textit{Bagaimana} sistem berperilaku. Contoh: load $<$ 3 detik, password di-hash, responsif. \textbf{Direkomendasikan}.
	\end{block}
	Identifikasi kebutuhan $\to$ tahap \textbf{analisis kebutuhan}; didokumentasikan dalam \textbf{SRS} (Software Requirements Specification).
\end{frame}

\begin{frame}{Perbandingan Kebutuhan Fungsional vs Non-Fungsional}
	\small
	\begin{tabular}{|l|l|l|}
		\hline
		\textbf{Aspek} & \textbf{Fungsional} & \textbf{Non-Fungsional} \\
		\hline
		Definisi & Apa yang dilakukan sistem & Bagaimana sistem berperilaku \\
		\hline
		Esensialitas & Wajib & Direkomendasikan \\
		\hline
		Hasil & Fitur aplikasi & Sifat/kualitas aplikasi \\
		\hline
		Fokus & Tuntutan pengguna & Harapan pengguna \\
		\hline
	\end{tabular}
\end{frame}

% ========== SECTION 2-3: Front-End ==========
\begin{frame}{Teknologi Front-End: HTML, CSS, JavaScript}
	\begin{itemize}
		\item \textbf{HTML}: struktur dan konten (kerangka rumah).
		\item \textbf{CSS}: presentasi visual (cat, dekorasi).
		\item \textbf{JavaScript}: interaktivitas (sistem listrik).
		\item Disebut \textit{holy trinity} pengembangan web.
	\end{itemize}
	\begin{block}{Peran}
		HTML = struktur \& konten; CSS = warna, font, layout; JS = validasi form, DOM, Fetch API.
	\end{block}
\end{frame}

\begin{frame}{HTML, CSS, JavaScript (lanjutan)}
	\begin{itemize}
		\item \textbf{HTML5}: elemen semantik (\texttt{<header>}, \texttt{<main>}, \texttt{<article>}), \texttt{<video>}, \texttt{<audio>}, \texttt{<canvas>}, tipe input baru.
		\item \textbf{CSS}: Box Model (content, padding, border, margin); \textbf{Flexbox} (satu dimensi); \textbf{Grid} (dua dimensi).
		\item \textbf{JavaScript}: DOM, event, ES6+ (let/const, arrow function, Promise, Fetch).
	\end{itemize}
	\begin{block}{Separation of Concerns}
		Struktur (HTML) $|$ Presentasi (CSS) $|$ Perilaku (JS). Simpan CSS dan JS di file terpisah.
	\end{block}
\end{frame}

% ========== SECTION 2-4: Back-End ==========
\begin{frame}{Teknologi Back-End}
	\begin{itemize}
		\item \textbf{PHP}: skrip sisi server; tag \texttt{<?php ... ?>}; dieksekusi di server, hasil (HTML) dikirim ke browser.
		\item \textbf{MySQL}: RDBMS; data dalam tabel; operasi \textbf{CRUD} dengan SQL (SELECT, INSERT, UPDATE, DELETE).
		\item \textbf{Web server}: Apache/Nginx menerima request, meneruskan ke PHP.
		\item \textbf{API/REST}: URL = sumber daya; metode HTTP = operasi. Format data: \textbf{JSON}.
	\end{itemize}
\end{frame}

\begin{frame}{Alur Request: Browser $\to$ Database}
	\begin{enumerate}
		\item DNS lookup (domain $\to$ IP)
		\item TCP handshake (port 443)
		\item HTTP Request (GET/POST, header)
		\item Web server $\to$ file PHP
		\item PHP $\to$ query SQL $\to$ MySQL
		\item MySQL $\to$ data $\to$ PHP
		\item PHP $\to$ HTML/JSON $\to$ HTTP Response
		\item Browser render (HTML + CSS + JS)
	\end{enumerate}
\end{frame}

% ========== SECTION 2-5: SDLC ==========
\begin{frame}{Siklus Pengembangan Web (SDLC)}
	\begin{enumerate}
		\item \textbf{Analisis kebutuhan}: kebutuhan fungsional \& non-fungsional $\to$ dokumen SRS.
		\item \textbf{Perancangan}: wireframe, mockup, user flow, arsitektur sistem.
		\item \textbf{Implementasi}: kode (HTML, CSS, JS, PHP); version control (Git).
		\item \textbf{Pengujian}: unit, integrasi, end-to-end; kompatibilitas browser.
		\item \textbf{Deployment}: server produksi, domain, DNS, SSL.
		\item \textbf{Pemeliharaan}: bug fix, update, monitoring, backup.
	\end{enumerate}
\end{frame}

\begin{frame}{Waterfall vs Agile}
	\begin{block}{Waterfall}
		Tahap berurutan dari awal sampai akhir. Cocok jika kebutuhan sudah jelas dan stabil.
	\end{block}
	\begin{block}{Agile}
		Iterasi pendek (\textbf{sprint}, 1--2 minggu); setiap sprint hasilkan fitur yang bisa diuji. Adaptif terhadap perubahan; populer di industri web.
	\end{block}
\end{frame}

% ========== SECTION 2-6: Alat ==========
\begin{frame}{Alat Pengembangan Web}
	\begin{itemize}
		\item \textbf{Editor}: VS Code (IntelliSense, terminal, Live Server, Prettier, PHP Intelephense, ESLint).
		\item \textbf{DevTools}: Elements (HTML/CSS), Console (JS), Network (request/response), Sources (debug).
		\item \textbf{Server lokal}: XAMPP (Apache + MySQL + PHP); proyek di \texttt{htdocs}; \texttt{http://localhost/}; phpMyAdmin.
		\item \textbf{Git}: \texttt{git init}, \texttt{git add}, \texttt{git commit}, \texttt{git push}; GitHub/GitLab.
		\item \textbf{Terminal}: perintah Git, server, npm/Composer.
	\end{itemize}
\end{frame}

% ========== SECTION 2-7: Standar Web ==========
\begin{frame}{Standar Web dan Organisasi}
	\begin{tabular}{|l|l|l|}
		\hline
		\textbf{Organisasi} & \textbf{Standar} & \textbf{Contoh} \\
		\hline
		W3C & CSS, SVG, ARIA & Flexbox, Grid, WCAG \\
		\hline
		WHATWG & HTML & HTML Living Standard \\
		\hline
		ECMA & JavaScript & ES2015 (ES6), ES2024 \\
		\hline
		IETF & Protokol & HTTP/1.1, HTTP/2, HTTP/3, TLS \\
		\hline
	\end{tabular}
	\vspace{0.5em}
	\begin{block}{Mengapa penting?}
		Interoperabilitas (konsisten di semua browser), aksesibilitas, masa depan kode. Validasi: \texttt{https://validator.w3.org/}
	\end{block}
\end{frame}

% ========== SECTION 2-8: Peta Konsep ==========
\begin{frame}{Peta Konsep: Hubungan Antar Teknologi}
	\begin{itemize}
		\item \textbf{Front-end}: Browser memuat HTML (struktur), CSS (tampilan), JavaScript (interaksi).
		\item \textbf{Back-end}: Web server $\to$ PHP $\to$ MySQL; response HTML/JSON.
		\item \textbf{Traditional}: setiap aksi = full page reload.
		\item \textbf{AJAX/Fetch}: JavaScript request data di latar belakang; hanya bagian yang berubah di-update; pengalaman lebih responsif.
	\end{itemize}
\end{frame}

% ========== SECTION 2-9: Keamanan ==========
\begin{frame}{Keamanan Web: Tiga Ancaman Utama}
	\begin{block}{SQL Injection}
		Input pengguna disisipkan ke query SQL. \textbf{Pencegahan}: \textbf{prepared statement}.
	\end{block}
	\begin{block}{XSS (Cross-Site Scripting)}
		Skrip berbahaya disisipkan ke halaman. \textbf{Pencegahan}: \textbf{output encoding} (\texttt{htmlspecialchars()}).
	\end{block}
	\begin{block}{CSRF (Cross-Site Request Forgery)}
		Request palsu memaksa aksi atas nama korban. \textbf{Pencegahan}: \textbf{CSRF token}, cookie \texttt{SameSite=Strict}.
	\end{block}
	\textbf{Prinsip}: \textit{Never trust user input}; validasi \& sanitasi di server.
\end{frame}

% ========== SECTION 2-10: Aksesibilitas dan Performa ==========
\begin{frame}{Aksesibilitas Web (a11y)}
	\begin{itemize}
		\item Web dapat digunakan semua orang (termasuk penyandang disabilitas).
		\item Standar \textbf{WCAG} (W3C WAI); prinsip \textbf{POUR}:
		\item \textbf{Perceivable}: konten dapat diakses (mis. teks alternatif gambar).
		\item \textbf{Operable}: antarmuka dapat dioperasikan (mis. keyboard).
		\item \textbf{Understandable}: informasi dan operasi dapat dipahami.
		\item \textbf{Robust}: konten dapat diinterpretasikan berbagai teknologi (HTML semantik, ARIA).
	\end{itemize}
\end{frame}

\begin{frame}{Performa: Core Web Vitals}
	\begin{itemize}
		\item \textbf{LCP} (Largest Contentful Paint): waktu hingga elemen terbesar selesai dirender. Target: $<$ 2,5 s.
		\item \textbf{INP} (Interaction to Next Paint): responsivitas interaksi pertama. Target: $<$ 200 ms.
		\item \textbf{CLS} (Cumulative Layout Shift): stabilitas visual (sedikit pergeseran). Target: $<$ 0,1.
		\item Optimasi: kompres gambar (WebP), \texttt{width/height} pada \texttt{<img>}, minifikasi CSS/JS, caching, lazy loading.
		\item \textbf{Lighthouse} (Chrome DevTools): audit Performance, Accessibility, Best Practices, SEO.
	\end{itemize}
\end{frame}

% ========== Rangkuman dan Penutup ==========
\begin{frame}{Rangkuman Bab 2}
	\begin{itemize}
		\item Kebutuhan fungsional vs non-fungsional (Sub-CPMK 1.1).
		\item Arsitektur client-server; HTTP; DNS; komponen URL; metode \& status code.
		\item Front-end: HTML, CSS, JavaScript; back-end: PHP, MySQL, API/REST.
		\item SDLC: analisis, perancangan, implementasi, pengujian, deployment, pemeliharaan; Agile vs Waterfall.
		\item Alat: VS Code, DevTools, XAMPP, Git.
		\item Standar: W3C, WHATWG, ECMA, IETF.
		\item Keamanan: SQL Injection, XSS, CSRF; defense in depth.
		\item Aksesibilitas (WCAG/POUR) dan performa (Core Web Vitals).
	\end{itemize}
\end{frame}

\begin{frame}{}
	\centering
	\vfill
	\textbf{\Large Pertanyaan?}
	\vfill
\end{frame}

\end{document}
